% problem_id: S001-lemma-K-injective-embedding-filtration
% source_id: S001 (Stacks Project)
% source_file: injectives.tex
% statement_locator: injectives.tex lines 2107--2125
% proof_locator: injectives.tex lines 2127--2240
% env: lemma
% id_note: label 在本来源内唯一
% pairing: tight (verified)
% status: complete (statement + author proof verbatim + reason + locators)
%
\begin{lemma}
\label{lemma-K-injective-embedding-filtration}
Let $\mathcal{A}$ be a Grothendieck abelian category.
Let $K^\bullet$ be a filtered complex of $\mathcal{A}$, see
Homology, Definition \ref{homology-definition-filtered-complex}.
Then there exists a morphism $j : K^\bullet \to J^\bullet$
of filtered complexes of $\mathcal{A}$ such that
\begin{enumerate}
\item $J^n$, $F^pJ^n$, $J^n/F^pJ^n$ and $F^pJ^n/F^{p'}J^n$ are injective
objects of $\mathcal{A}$,
\item $J^\bullet$, $F^pJ^\bullet$, $J^\bullet/F^pJ^\bullet$, and
$F^pJ^\bullet/F^{p'}J^\bullet$ are K-injective complexes,
\item $j$ induces quasi-isomorphisms
$K^\bullet \to J^\bullet$,
$F^pK^\bullet \to F^pJ^\bullet$,
$K^\bullet/F^pK^\bullet \to J^\bullet/F^pJ^\bullet$, and
$F^pK^\bullet/F^{p'}K^\bullet \to F^pJ^\bullet/F^{p'}J^\bullet$.
\end{enumerate}
\end{lemma}

\begin{proof}
By Theorem \ref{theorem-K-injective-embedding-grothendieck}
we obtain quasi-isomorphisms $i : K^\bullet \to I^\bullet$ and
$i^p : F^pK^\bullet \to I^{p, \bullet}$ as well as commutative diagrams
$$
\vcenter{
\xymatrix{
K^\bullet \ar[d]_i & F^pK^\bullet \ar[l] \ar[d]_{i^p} \\
I^\bullet & I^{p, \bullet} \ar[l]_{\alpha^p}
}
}
\quad\text{and}\quad
\vcenter{
\xymatrix{
F^{p'}K^\bullet \ar[d]_{i^{p'}} &
F^pK^\bullet \ar[l] \ar[d]_{i^p} \\
I^{p', \bullet} &
I^{p, \bullet} \ar[l]_{\alpha^{p p'}}
}
}
\quad\text{for }p' \leq p
$$
such that $\alpha^p \circ \alpha^{p' p} = \alpha^{p'}$
and $\alpha^{p'p''} \circ \alpha^{pp'} = \alpha^{pp''}$.
The problem is that the maps $\alpha^p : I^{p, \bullet} \to I^\bullet$
need not be injective. For each $p$ we choose an injection
$t^p : I^{p, \bullet} \to J^{p, \bullet}$ into an acyclic K-injective
complex $J^{p, \bullet}$ whose terms are injective objects of $\mathcal{A}$
(first map to the cone on the identity and then use the theorem).
Choose a map of complexes $s^p : I^\bullet \to J^{p, \bullet}$
such that the following diagram commutes
$$
\xymatrix{
K^\bullet \ar[d]_i & F^pK^\bullet \ar[l] \ar[d]_{i^p} \\
I^\bullet \ar[rd]_{s^p} & I^{p, \bullet} \ar[d]^{t^p} \\
& J^{p, \bullet}
}
$$
This is possible: the composition $F^pK^\bullet \to J^{p, \bullet}$
is homotopic to zero because $J^{p, \bullet}$ is acyclic and K-injective
(Derived Categories, Lemma \ref{derived-lemma-K-injective}).
Since the objects $J^{p, n - 1}$ are injective and since
$F^pK^n \to K^n \to I^n$ are injective morphisms, we
can lift the maps $F^pK^n \to J^{p, n - 1}$ giving the homotopy
to a map $h^n : I^n \to J^{p, n - 1}$. Then we set $s^p$
equal to $h \circ \text{d} + \text{d} \circ h$.
(Warning: It will not be the case that $t^p = s^p \circ \alpha^p$,
so we have to be careful not to use this below.)

\medskip\noindent
Consider
$$
J^\bullet = I^\bullet \times \prod\nolimits_p J^{p, \bullet}
$$
Because products in $D(\mathcal{A})$ are given by taking
products of K-injective complexes
(Lemma \ref{lemma-derived-products})
and since $J^{p, \bullet}$
is isomorphic to $0$ in $D(\mathcal{A})$ we see that
$J^\bullet \to I^\bullet$ is an isomorphism in $D(\mathcal{A})$.
Consider the map
$$
j = i \times (s^p \circ i)_{p \in \mathbf{Z}} :
K^\bullet
\longrightarrow
I^\bullet \times \prod\nolimits_p J^{p, \bullet} = J^\bullet
$$
By our remarks above this is a quasi-isomorphism. It is also injective.
For $p \in \mathbf{Z}$ we let $F^pJ^\bullet \subset J^\bullet$ be
$$
\Im\left(
\alpha^p \times (t^{p'} \circ \alpha^{pp'})_{p' \leq p} :
I^{p, \bullet}
\to
I^\bullet \times \prod\nolimits_{p' \leq p} J^{p', \bullet}
\right)
\times \prod\nolimits_{p' > p} J^{p', \bullet}
$$
This complex is isomorphic to the complex
$I^{p, \bullet} \times \prod_{p' > p} J^{p, \bullet}$
as $\alpha^{pp} = \text{id}$ and $t^p$ is injective.
Hence $F^pJ^\bullet$ is quasi-isomorphic to $I^{p, \bullet}$ (argue
as above). We have $j(F^pK^\bullet) \subset F^pJ^\bullet$ because
of the commutativity of the diagram above. The corresponding
map of complexes $F^pK^\bullet \to F^pJ^\bullet$ is a quasi-isomorphism
by what we just said. Finally, to see that
$F^{p + 1}J^\bullet \subset F^pJ^\bullet$
use that $\alpha^{p + 1p} \circ \alpha^{pp'} = \alpha^{p + 1p'}$
and the commutativity of the first displayed diagram
in the first paragraph of the proof.

\medskip\noindent
We claim that $j : K^\bullet \to J^\bullet$ is a solution to the
problem posed by the lemma. Namely, $F^pJ^n$ is an injective object
of $\mathcal{A}$ because it is isomorphic to
$I^{p, n} \times \prod_{p' > p} J^{p', n}$ and products of
injectives are injective. Then the injective map $F^pJ^n \to J^n$
splits and hence the quotient $J^n/F^pJ^n$ is injective as well
as a direct summand of the injective object $J^n$.
Similarly for $F^pJ^n/F^{p'}J^n$. This in particular means
that $0 \to F^pJ^\bullet \to J^\bullet \to J^\bullet/F^pJ^\bullet \to 0$
is a termwise split short exact sequence of complexes, hence defines
a distinguished triangle in $K(\mathcal{A})$ by fiat.
Since $J^\bullet$ and $F^pJ^\bullet$ are K-injective complexes
we see that the same is true for $J^\bullet/F^pJ^\bullet$
by Derived Categories, Lemma \ref{derived-lemma-triangle-K-injective}.
A similar argument shows that $F^pJ^\bullet/F^{p'}J^\bullet$
is K-injective. By construction $j : K^\bullet \to J^\bullet$
and the induced maps $F^pK^\bullet \to F^pJ^\bullet$ are
quasi-isomorphisms. Using the long exact cohomology sequences
of the complexes in play we find that the same holds for
$K^\bullet/F^pK^\bullet \to J^\bullet/F^pJ^\bullet$ and
$F^pK^\bullet/F^{p'}K^\bullet \to F^pJ^\bullet/F^{p'}J^\bullet$.
\end{proof}
